\documentclass{article}
\usepackage{amsmath,amssymb,amsthm,mathtools}
\usepackage{proofread}
\newtheorem{proposition}{Proposition}
\newcommand{\CC}{\mathcal{C}}

\title{Tail monotone copulas are positively quadrant dependent}
\author{Proofread example}
\date{}

\begin{document}
\maketitle

\begin{abstract}
A copula is left-tail decreasing when the conditional probability of a small second coordinate decreases in the first.
We show in a few lines that such copulas, and likewise right-tail increasing ones, are positively quadrant dependent.
\end{abstract}

\section{Introduction}

Positive dependence of two random variables can be expressed by several orderings of their copula, and \cite[Ch.~5]{nelsen2006introduction} collects the classical ones.
Tail monotonicity is among the strongest of them.
The implications between these notions are well known, and we include the short arguments for the two tail notions below.

\section{Notions of positive dependence}

Let $\CC$ denote the class of bivariate copulas.
For the copula-theoretic formulations used here, we follow \cite[Ch.~5]{nelsen2006introduction}.
We call a copula $C$ \emph{positively quadrant dependent} (\emph{PQD}) if $C(u,v)\ge uv$ for all $u,v\in[0,1]$.
A copula $C$ is called \emph{left-tail decreasing} (\emph{LTD}) if, for every $v\in[0,1]$, the map
\begin{equation}\label{eq:ltd}
u\longmapsto\frac{C(u,v)}{u},\qquad u\in(0,1],
\end{equation}
is non-increasing.
Similarly, $C$ is called \emph{right-tail increasing} (\emph{RTI}) if, for every $v\in[0,1]$, the map
\begin{equation}\label{eq:rti}
u\longmapsto\frac{1-u-v+C(u,v)}{1-u},\qquad u\in[0,1),
\end{equation}
is non-decreasing.

\section{Tail monotonicity implies quadrant dependence}

\begin{proposition}\label{prop:ltd}
Every LTD copula is PQD.
\end{proposition}

\begin{proof}
Let $C$ be an LTD copula, and let $u,v\in[0,1]$.
If $u=0$, then $C(u,v)=0=uv$.
If $u>0$, then $\frac{C(u,v)}{u}\ge C(1,v)=v$ by \eqref{eq:ltd}.
Hence $C(u,v)\ge uv$.
\end{proof}

\begin{proposition}\label{prop:rti}
Every RTI copula is PQD.
\end{proposition}

\begin{proof}
Let $C$ be an RTI copula, and let $u,v\in[0,1]$.
If $u=1$, then $C(u,v)=v=uv$.
If $u<1$, then $\frac{1-u-v+C(u,v)}{1-u}\ge1-v+C(0,v)=1-v$ by \eqref{eq:rti}.
Hence $C(u,v)\ge uv$.
\end{proof}

\begin{thebibliography}{1}
\bibitem{nelsen2006introduction} R.~B. Nelsen, \emph{An Introduction to Copulas}, 2nd ed., Springer Series in Statistics, Springer, New York, 2006.
\end{thebibliography}

\end{document}
